% =============================================================================
%  Notas de clase — Sesión 22: Archivos binarios
%  Programación Avanzada — Universidad Panamericana
%  Compilar: pdflatex sesion22_archivos_binarios.tex
% =============================================================================
\documentclass[12pt, oneside]{article}
\usepackage{progavanzada}
\usepackage{array}

\begin{document}

\portada{Sesión 22}{Archivos binarios}{2026}

\tableofcontents
\newpage

% =============================================================================
\section{1973: los bytes que se reconocen solos}
% =============================================================================

En la sesión anterior viste que un archivo no es más que una secuencia de
bytes almacenada en disco.  Esa definición aplica tanto a un \texttt{.txt}
como a un \texttt{.jpg}, un \texttt{.mp3} o el ejecutable \texttt{a.out}
que produces cada vez que compilas tu programa.  La diferencia no está en
el soporte físico sino en \textit{cómo se interpretan} esos bytes.

En 1973, los laboratorios Bell publicaron la Cuarta Edición de Unix.  Entre
las herramientas que incluyó había una pequeña pero ingeniosa: el comando
\textbf{\texttt{file}}.  Su función era leer los primeros bytes de un
archivo y determinar de qué tipo era, sin depender del nombre ni de la
extensión.  La herramienta era posible porque los diseñadores de formatos
habían adoptado una convención: los primeros bytes de cada formato forman
una \textbf{firma} (o \textbf{número mágico}) que lo identifica de forma
única.

Hoy esa convención es universal.  Cuando abres una foto, escuchas una
canción o ejecutas un programa, el sistema operativo lee los primeros bytes
del archivo para saber cómo procesarlo.

\begin{center}
\begin{tabular}{llll}
  \toprule
  \textbf{Formato} & \textbf{Firma (hex)}            & \textbf{ASCII}          & \textbf{Uso} \\
  \midrule
  ELF              & \texttt{7F 45 4C 46}             & \texttt{\textbackslash x7FELF} & Ejecutables Linux \\
  BMP              & \texttt{42 4D}                   & \texttt{BM}             & Imágenes (Windows) \\
  JPEG             & \texttt{FF D8 FF}                & —                       & Imágenes \\
  PNG              & \texttt{89 50 4E 47}             & \texttt{\textbackslash x89PNG} & Imágenes \\
  MP3 (ID3)        & \texttt{49 44 33}                & \texttt{ID3}            & Audio \\
  WAV              & \texttt{52 49 46 46}             & \texttt{RIFF}           & Audio \\
  ZIP              & \texttt{50 4B 03 04}             & \texttt{PK}             & Archivos comprimidos \\
  PDF              & \texttt{25 50 44 46}             & \texttt{\%PDF}          & Documentos \\
  \bottomrule
\end{tabular}
\end{center}

Aquí viene la sorpresa: el archivo \texttt{a.out} que produces cuando
compilas con \texttt{g++} es un archivo binario en formato \textbf{ELF}
(\textit{Executable and Linkable Format}).  Sus primeros cuatro bytes son
siempre \texttt{7F 45 4C 46}.  Puedes comprobarlo en tu terminal:

\begin{verbatim}
$ g++ hola.cpp -o hola
$ file hola
hola: ELF 64-bit LSB pie executable, x86-64, ...
$ xxd hola | head -1
00000000: 7f45 4c46 0201 0100 0000 0000 0000 0000  .ELF............
\end{verbatim}

Llevas un semestre produciendo archivos binarios sin saberlo.  Hoy vas a
entender cómo funcionan y cómo crearlos tú mismo.

\begin{tecnico}[El comando \texttt{file} y los números mágicos]
  El término ``número mágico'' viene del código fuente original de Unix.
  Los primeros dos bytes de los ejecutables PDP-11 eran el valor octal
  \texttt{0407}, que el \textit{loader} buscaba antes de ejecutar el
  archivo.  Hoy la base de datos de firmas conocidas
  (\texttt{/usr/share/misc/magic}) contiene miles de entradas, desde
  formatos de los años~70 hasta los más recientes.
  \par\smallskip
  \textit{Referencia:} Ritchie, D.\,M. (1984). ``The Evolution of the
  Unix Time-sharing System.'' \textit{AT\&T Bell Laboratories Technical
  Journal}, 63(8), 1577--1593.
\end{tecnico}

% =============================================================================
\section{El problema del modo texto}
% =============================================================================

En la sesión anterior guardaste enteros en archivos de texto con el
operador \cpp{<<}.  Funciona bien para mostrar resultados, pero tiene un
problema cuando los datos deben ser leídos por otro programa.

Considera el siguiente código:

\textbf{binarios1.cpp}:
\begin{verbatim}
#include <iostream>
#include <fstream>
using namespace std;

int main(){
    int x = 123;
    int y = 27;

    ofstream f1("datos.out");
    f1 << x << y;
    f1.close();

    ifstream f2("datos.out");
    int z;
    f2 >> z;
    cout << z << endl;
}
\end{verbatim}

¿Qué imprime?  Antes de leer la respuesta, razona: ¿qué contiene
\texttt{datos.out} después de la escritura?

El operador \cpp{<<} convierte los enteros a su representación en texto.
\texttt{123} se convierte en los tres caracteres \texttt{'1'}, \texttt{'2'},
\texttt{'3'}, y \texttt{27} en los dos caracteres \texttt{'2'},
\texttt{'7'}.  El archivo queda así:

\begin{center}
\begin{tabular}{ccccc}
  \toprule
  Byte 0 & Byte 1 & Byte 2 & Byte 3 & Byte 4 \\
  \midrule
  \texttt{0x31} & \texttt{0x32} & \texttt{0x33} & \texttt{0x32} & \texttt{0x37} \\
  \texttt{'1'}  & \texttt{'2'}  & \texttt{'3'}  & \texttt{'2'}  & \texttt{'7'} \\
  \bottomrule
\end{tabular}
\end{center}

Cinco bytes que forman la cadena \texttt{"12327"}.  Cuando \cpp{f2 >> z}
intenta leer un entero, lee todos los dígitos hasta el final del archivo y
obtiene \texttt{12327}.  El programa imprime \textbf{12327}, no \textbf{123}.

La solución parcial en modo texto es agregar un separador:
\begin{verbatim}
f1 << x << " " << y;
\end{verbatim}

Pero esto tiene dos problemas:

\begin{enumerate}
  \item Hay que recordar poner el separador en la escritura y manejarlo en
        la lectura.  Si los datos son más complejos (cadenas con espacios,
        tipos mixtos), el parseo se vuelve frágil.
  \item Es ineficiente: el número \texttt{2147483647} (el mayor \cpp{int})
        ocupa 10 bytes como texto, pero en binario siempre ocupa 4.
\end{enumerate}

El \textbf{modo binario} resuelve ambos problemas de raíz.

\begin{recuerda}
  ``Binario'' no significa que el archivo contenga ceros y unos escritos
  como texto.  Todos los archivos son binarios a nivel del disco.
  ``Modo binario'' significa que no hay conversión entre el dato en memoria
  y los bytes en disco: se copian los bytes exactos tal como están en la
  RAM.
\end{recuerda}

% =============================================================================
\section{El modo binario: \texttt{write} y \texttt{read}}
% =============================================================================

Para abrir un archivo en modo binario, agrega la bandera
\cpp{fstream::binary} al constructor.  Luego, en lugar de \cpp{<<} y
\cpp{>>}, se usan los métodos \cpp{write} y \cpp{read}:

\textbf{binarios2.cpp}:
\begin{verbatim}
#include <iostream>
#include <fstream>
using namespace std;

int main(){
    int x = 123;
    int y = 27;

    ofstream f1("datosbin.out", fstream::binary);
    f1.write( (char *) &x, sizeof(int) );
    f1.write( (char *) &y, sizeof(int) );
    f1.close();

    ifstream f2("datosbin.out", fstream::binary);
    int w, z;
    f2.read( (char *) &w, sizeof(int) );
    f2.read( (char *) &z, sizeof(int) );
    cout << w << " " << z << endl;
}
\end{verbatim}

Salida:
\begin{verbatim}
123 27
\end{verbatim}

El archivo \texttt{datosbin.out} tiene exactamente \textbf{8 bytes}:
4 para \texttt{x} y 4 para \texttt{y}, copiados directamente de la memoria.

\subsection{Desglose de \texttt{write} y \texttt{read}}

Ambos métodos tienen la misma firma:
\begin{verbatim}
ostream& write(const char* ptr, streamsize n);
istream& read(char* ptr, streamsize n);
\end{verbatim}

\begin{description}
  \item[\cpp{ptr}]
    Puntero a los bytes que se van a escribir (o donde se van a colocar
    los bytes leídos).  El tipo es \cpp{const char*}, que en C++ es el
    tipo básico de ``puntero a bytes''.

  \item[\cpp{n}]
    Número de bytes a escribir o leer.  Casi siempre es \cpp{sizeof}
    del tipo que quieres guardar.
\end{description}

\subsection{El cast \texttt{(char*) \&x}}

Esta es la parte que más confunde.  Analízala en dos partes:

\begin{enumerate}
  \item \cpp{\&x} es la dirección de memoria de la variable \cpp{x}.
        Su tipo es \cpp{int*} (``puntero a int'').
  \item \cpp{(char *)} es un \textbf{cast}: le dices al compilador
        ``trata esta dirección como un puntero a bytes, no como un
        puntero a int''.
\end{enumerate}

No estás cambiando los datos.  La variable \cpp{x} sigue siendo un
\cpp{int} con el valor 123.  Solo le estás diciendo a \cpp{write}:
``empieza en esta dirección y copia los próximos \cpp{sizeof(int)} bytes
tal como están''.

\begin{advertencia}
  El cast \cpp{(char*)} es necesario porque \cpp{write} no sabe si los
  datos son \cpp{int}, \cpp{double}, un struct u otra cosa.  Solo trabaja
  con bytes.  El programador es responsable de indicar la dirección
  correcta y el tamaño correcto.  Si el tamaño que pasas a \cpp{write}
  no coincide con el tamaño real del dato, los datos llegan corruptos
  al disco, y el error suele ser silencioso.
\end{advertencia}

\begin{ejemplo}[Verificar el tamaño del archivo]
  Después de ejecutar \textbf{binarios2.cpp} puedes confirmar que el
  archivo tiene exactamente 8 bytes:
  \begin{verbatim}
$ ls -l datosbin.out
-rw-r--r-- 1 usuario grupo 8 may 11 datosbin.out
  \end{verbatim}
  Compara con el archivo en modo texto: tendría 5 bytes (\texttt{"12327"})
  o 6 con el separador (\texttt{"123 27"}).  El archivo binario es más
  pequeño y no necesita parseo.
\end{ejemplo}

% =============================================================================
\section{Guardar structs completos}
% =============================================================================

El verdadero poder del modo binario aparece cuando trabajas con
\textbf{structs}.  Un struct es un bloque contiguo de bytes en memoria.
\cpp{write} puede copiar ese bloque completo al disco en una sola llamada,
sin necesidad de escribir campo por campo.

Reutilizamos el struct de la actividad anterior:

\textbf{persona.h}:
\begin{verbatim}
using namespace std;

struct Persona {
    char nombre[30];
    char apellido[30];
    int  dia;
    int  mes;
};
\end{verbatim}

\subsection{Escribir un struct}

\begin{verbatim}
#include <iostream>
#include <fstream>
#include <cstring>
#include "persona.h"
using namespace std;

int main(){
    Persona p;
    strcpy(p.nombre,   "Hannah");
    strcpy(p.apellido, "Fabián");
    p.dia = 29;
    p.mes = 2;

    ofstream f("hannah.cumple", fstream::binary);
    f.write( (char *) &p, sizeof(Persona) );
    f.close();
}
\end{verbatim}

\cpp{sizeof(Persona)} calcula el tamaño del struct en bytes: en principio
$30 + 30 + 4 + 4 = 68$, aunque el compilador puede agregar relleno
(\textit{padding}); ver el ejercicio~3.

\subsection{Leer un struct}

\begin{verbatim}
#include <iostream>
#include <fstream>
#include "persona.h"
using namespace std;

int main(){
    Persona p;

    ifstream f("hannah.cumple", fstream::binary);
    f.read( (char *) &p, sizeof(Persona) );
    f.close();

    cout << p.nombre  << " " << p.apellido
         << ": " << p.dia << "/" << p.mes << endl;
}
\end{verbatim}

Salida:
\begin{verbatim}
Hannah Fabián: 29/2
\end{verbatim}

\begin{consejo}
  La simetría es total: escritura y lectura usan \cpp{\&p} y
  \cpp{sizeof(Persona)}.  Si en alguno de los dos lados cambias el
  tamaño o la dirección, los datos llegan corruptos.  Escribe primero,
  lee inmediatamente y verifica la salida antes de continuar.
\end{consejo}

\subsection{Múltiples registros y acceso aleatorio}

Para guardar varios registros, llama a \cpp{write} una vez por registro:

\begin{verbatim}
ofstream f("agenda.bin", fstream::binary);
for (int i = 0; i < n; i++){
    f.write( (char *) &personas[i], sizeof(Persona) );
}
f.close();
\end{verbatim}

El archivo resultante tiene exactamente \cpp{n * sizeof(Persona)} bytes.
Para leer directamente el registro en la posición~$k$ (sin leer los
anteriores), usa \cpp{seekg} para saltar:

\begin{verbatim}
ifstream f("agenda.bin", fstream::binary);
Persona p;
f.seekg(k * sizeof(Persona));
f.read( (char *) &p, sizeof(Persona) );
\end{verbatim}

Esto se llama \textbf{acceso aleatorio}: puedes ir directo a cualquier
registro como si fuera el elemento $k$ de un arreglo.  Con archivos de
texto no es posible hacerlo de forma confiable porque los registros tienen
longitud variable.

\begin{tecnico}[Por qué arreglos de \texttt{char} en lugar de \texttt{string}]
  El struct \cpp{Persona} usa \cpp{char nombre[30]} en lugar de
  \cpp{string}.  Un \cpp{string} en C++ no es un bloque contiguo de
  bytes: es un objeto que guarda internamente un puntero a memoria
  dinámica.  Si lo escribieras en binario, guardarías la dirección
  del puntero, no el texto.  En otra ejecución esa dirección no
  tendría sentido.  Los arreglos de \cpp{char} son simples bloques
  de bytes y por eso funcionan correctamente con \cpp{write} y
  \cpp{read}.
\end{tecnico}

% =============================================================================
\section{Formatos binarios en el mundo real}
% =============================================================================

Todos los formatos de archivo que usas a diario son binarios: imágenes,
canciones, documentos, ejecutables.  Cada uno tiene una estructura definida
por un estándar y es, en esencia, un conjunto de structs escritos en disco.
A continuación verás cómo están construidos algunos de los más comunes.

\subsection{BMP: el formato de imagen más simple}

El formato \textbf{BMP} (\textit{BitMaP}) fue introducido con Windows~3.0
en 1990.  No aplica compresión: los píxeles se almacenan directamente,
lo que lo hace simple de leer pero voluminoso.  Una imagen de
$100\times100$ píxeles a 24~bits ocupa $100\times100\times3=30{,}000$
bytes de datos de imagen, más el encabezado.

La estructura es casi directamente expresable como structs de C++:

\begin{verbatim}
// Encabezado del archivo (14 bytes)
struct BMPHeader {
    char  firma[2];      // siempre "BM"
    int   tamano;        // tamaño total del archivo en bytes
    short reservado1;    // siempre 0
    short reservado2;    // siempre 0
    int   offset_datos;  // posición donde empiezan los píxeles
};

// Encabezado de imagen (40 bytes)
struct BMPInfo {
    int    tam_encabezado; // siempre 40
    int    ancho;          // píxeles de ancho
    int    alto;           // píxeles de alto
    short  planos;         // siempre 1
    short  bits_pixel;     // 24 para color completo
    int    compresion;     // 0 = sin compresión
    int    tam_imagen;     // tamaño de los datos (puede ser 0)
    int    ppm_x;          // píxeles por metro horizontal
    int    ppm_y;          // píxeles por metro vertical
    int    colores_tabla;  // 0 = todos
    int    colores_imp;    // 0 = todos
};
\end{verbatim}

Después de los 54 bytes de encabezado vienen los píxeles en orden
\textbf{BGR} (azul, verde, rojo), de la fila inferior a la superior.
El orden invertido es un legado de la arquitectura \textit{little-endian}
de Intel.

\subsection{JPEG: compresión con pérdida}

El formato \textbf{JPEG} fue estandarizado en 1992 por el
\textit{Joint Photographic Experts Group}.  A diferencia de BMP, aplica
compresión con pérdida: descarta información que el ojo humano percibe
poco (frecuencias altas en el espacio de color YCbCr) y reduce el tamaño
del archivo típicamente entre 10:1 y 20:1.  Sin JPEG, las páginas web con
fotografías habrían sido imprácticamente lentas en la era del módem.

Un JPEG no almacena píxeles directamente.  Sus datos se organizan en
\textbf{segmentos}, cada uno precedido por un marcador de dos bytes que
comienza con \texttt{FF}:

\begin{center}
\begin{tabular}{lll}
  \toprule
  \textbf{Marcador} & \textbf{Hex} & \textbf{Significado} \\
  \midrule
  SOI  & \texttt{FF D8} & Inicio del archivo (\textit{Start Of Image}) \\
  APP0 & \texttt{FF E0} & Información JFIF (versión, resolución) \\
  SOF0 & \texttt{FF C0} & Parámetros de imagen (ancho, alto, canales) \\
  DHT  & \texttt{FF C4} & Tablas Huffman (para descomprimir) \\
  SOS  & \texttt{FF DA} & Inicio de los datos comprimidos \\
  EOI  & \texttt{FF D9} & Fin del archivo (\textit{End Of Image}) \\
  \bottomrule
\end{tabular}
\end{center}

Por eso los primeros dos bytes de todo JPEG son siempre \texttt{FF D8}:
marcan el inicio del archivo.  El último marcador que encontrarás al
final es siempre \texttt{FF D9}.

\subsection{WAV: audio sin compresión}

El formato \textbf{WAV} guarda audio en formato PCM
(\textit{Pulse-Code Modulation}): muestras numéricas tomadas a intervalos
regulares.  Un CD de audio usa 44\,100 muestras por segundo, cada una de
16~bits.  Un minuto en estéreo son
$44{,}100 \times 2 \times 2 \times 60 \approx 10$~MB.

WAV usa el contenedor \textbf{RIFF} (\textit{Resource Interchange File
Format}), una secuencia de bloques donde cada uno tiene un identificador
de cuatro caracteres, un tamaño y sus datos.  El inicio es siempre:

\begin{verbatim}
"RIFF"  [tamaño total - 8 bytes]  "WAVE"
"fmt "  [16]  [1=PCM]  [canales]  [frecuencia_Hz]  [...]
"data"  [tamaño de las muestras]  [muestras PCM...]
\end{verbatim}

El espacio en \texttt{"fmt "} es parte del identificador y no puede
omitirse.  Las muestras son simplemente enteros de 16 o 24~bits escritos
en modo binario, uno tras otro.

\subsection{ELF: el ejecutable que ya conoces}

El formato \textbf{ELF} (\textit{Executable and Linkable Format}) es el
estándar para ejecutables y bibliotecas en Linux.  Cada \texttt{a.out}
que has producido este semestre es un ELF.  Su encabezado inicial ocupa
64~bytes en arquitecturas de 64~bits:

\begin{verbatim}
struct ELFHeader {
    unsigned char magic[4];  // 7F 45 4C 46 = "\x7FELF"
    unsigned char clase;     // 1=32 bits, 2=64 bits
    unsigned char endian;    // 1=little-endian, 2=big-endian
    unsigned char version;   // siempre 1
    // ... campos de relleno ...
    unsigned short tipo;     // 2=ejecutable, 3=librería compartida
    unsigned short maquina;  // 62=x86-64
    // ...
};
\end{verbatim}

Después del encabezado vienen las secciones: \texttt{.text} (código
máquina), \texttt{.data} (variables globales inicializadas), \texttt{.bss}
(variables no inicializadas) y otras que el enlazador y el cargador del
sistema operativo usan para ejecutar el programa.

\begin{consejo}
  Los formatos binarios son, en esencia, structs escritos en disco con
  una estructura definida por un estándar.  Las habilidades que aprendiste
  hoy son suficientes para leer la cabecera de un BMP, verificar los bytes
  mágicos de cualquier archivo, o diseñar tus propios formatos.  Las bases
  de datos de alto rendimiento como SQLite usan exactamente este principio:
  páginas de tamaño fijo con una estructura bien definida.
\end{consejo}

% =============================================================================
\section{Actividad: agenda de cumpleaños binaria}
% =============================================================================

Esta actividad es la continuación directa de la sesión~21, pero ahora
los datos se guardan y leen en formato binario usando el struct
\cpp{Persona}.

\subsection{Formato del archivo}

Cada archivo personal es un único registro binario de tipo \cpp{Persona}
(\cpp{sizeof(Persona)} bytes).  El nombre del archivo es tu nombre en
minúsculas con extensión \texttt{.cumple} (por ejemplo,
\texttt{hannah.cumple}).

\subsection{Parte A: genera tu archivo (10 min)}

Escribe un programa que cree tu propio archivo \texttt{.cumple}:

\begin{enumerate}
  \item Define una variable de tipo \cpp{Persona} y llénala con tus datos
        reales usando \cpp{strcpy}.
  \item Ábrela en modo binario con \cpp{ofstream} y escribe el struct.
  \item Cierra el archivo.
  \item Abre el mismo archivo con \cpp{ifstream} en modo binario, lee el
        struct de vuelta y muestra los datos en pantalla para verificar
        que la escritura fue correcta.
\end{enumerate}

Sube tu archivo \texttt{.cumple} a la carpeta de Google Drive indicada por
el profesor.

\subsection{Parte B: las consultas (40 min)}

Descarga todos los archivos \texttt{.cumple} de tus compañeros y
colócalos en una carpeta.  Luego resuelve los siguientes problemas.
Para cada uno escribe un programa separado.

\begin{enumerate}

  \item \textbf{Unir todos los registros en un solo archivo binario.}\\
        Lee cada archivo \texttt{.cumple} individual y escribe todos los
        registros en un archivo llamado \texttt{agenda.bin}.  El resultado
        es un arreglo de \cpp{Persona} en disco: cada registro ocupa
        exactamente \cpp{sizeof(Persona)} bytes, uno tras otro.  Puedes
        apoyarte en un archivo de texto auxiliar con los nombres de los
        archivos individuales, uno por línea.

  \item \textbf{¿Cuántas personas hay en la agenda?}\\
        Abre \texttt{agenda.bin}, mueve el stream al final con
        \cpp{seekg(0, ios::end)} y consulta la posición con \cpp{tellg()}.
        Divide entre \cpp{sizeof(Persona)}.  Muestra el resultado.

  \item \textbf{Buscar por nombre.}\\
        Pide un nombre al usuario y recorre \texttt{agenda.bin} registro
        por registro hasta encontrar una coincidencia.  Si existe, muestra
        el nombre completo y la fecha de cumpleaños.

  \item \textbf{Acceso aleatorio: el registro $k$.}\\
        Pide un número $k$ al usuario y muestra directamente el
        $k$-ésimo registro de \texttt{agenda.bin} usando \cpp{seekg},
        sin leer los registros anteriores.

\end{enumerate}

% =============================================================================
\section{Ejercicios}
% =============================================================================

\begin{enumerate}

  \item Modifica \textbf{binarios2.cpp} para guardar y leer tres valores
        de tipo \cpp{double} en lugar de dos \cpp{int}.  Verifica con
        \texttt{ls -l} que el archivo resultante tiene exactamente
        \cpp{3 * sizeof(double)} bytes.

  \item Define un struct \cpp{Producto} con los campos:
        \cpp{char nombre[40]}, \cpp{double precio} e \cpp{int cantidad}.
        Escribe un programa que guarde cinco productos en un archivo
        binario llamado \texttt{inventario.bin} y luego los lea todos y
        los muestre en pantalla.

  \item \cpp{sizeof(Persona)} no siempre es $30 + 30 + 4 + 4 = 68$.
        Compila y ejecuta:
        \begin{verbatim}
#include <iostream>
#include "persona.h"
using namespace std;
int main(){
    cout << sizeof(Persona) << endl;
}
        \end{verbatim}
        Si el resultado es mayor que~68, el compilador está agregando
        \textbf{relleno} (\textit{padding}) para alinear los campos en
        memoria.  ¿Crees que esto afecta la escritura y lectura binaria
        si el escritor y el lector compilan con el mismo compilador?
        ¿Y si los archivos se comparten entre sistemas distintos?

  \item Escribe un programa que abra cualquier archivo (cuyo nombre pide
        al usuario), lea sus primeros 8 bytes e imprima cada byte en
        hexadecimal con el formato \texttt{0xNN}.  Pruébalo con un
        \texttt{.jpg}, un \texttt{.png} y el ejecutable \texttt{a.out}.
        Compara los resultados con la tabla de números mágicos de la
        sección~1.

  \item Extiende el ejercicio~2 para modificar el precio de un producto
        dado su índice (0 a~4) \textit{sin reescribir todo el archivo}.
        Usa \cpp{fstream} con los modos
        \cpp{fstream::in | fstream::out | fstream::binary} y el método
        \cpp{seekp} para posicionarte en el registro correcto antes de
        llamar a \cpp{write}.

\end{enumerate}

% =============================================================================
\section*{Resumen}
% =============================================================================

\begin{itemize}
  \item En \textbf{modo texto} los datos se convierten a caracteres al
        escribir y se parsean al leer; en \textbf{modo binario} los bytes
        se copian directamente entre la memoria y el disco.
  \item La bandera \cpp{fstream::binary} activa el modo binario al abrir
        el archivo.
  \item \cpp{write((char*)\&dato, sizeof(tipo))} escribe los bytes exactos
        del dato; \cpp{read((char*)\&dato, sizeof(tipo))} los recupera.
  \item El cast \cpp{(char*)\&x} indica a \cpp{write}/\cpp{read} dónde
        están los bytes; \cpp{sizeof} indica cuántos copiar.  Son dos
        informaciones distintas y ambas deben ser correctas.
  \item Los structs con campos de tamaño fijo (\cpp{char[n]}, \cpp{int},
        \cpp{double}) se serializan directamente con \cpp{write}/\cpp{read}.
        Los \cpp{string} no, porque internamente contienen punteros.
  \item El acceso aleatorio a registros de tamaño fijo se logra con
        \cpp{seekg(k * sizeof(T))} para lectura y
        \cpp{seekp(k * sizeof(T))} para escritura.
  \item Todos los formatos de archivo populares (JPEG, PNG, WAV, ELF, BMP)
        son structs escritos en disco con una estructura definida por un
        estándar.
  \item Los primeros bytes de cada formato forman una \textbf{firma} o
        \textbf{número mágico} que identifica el tipo de archivo sin
        depender del nombre ni de la extensión.
\end{itemize}

\end{document}
